% Options for packages loaded elsewhere
\PassOptionsToPackage{unicode}{hyperref}
\PassOptionsToPackage{hyphens}{url}
\documentclass[
  ignorenonframetext,
]{beamer}
\newif\ifbibliography
\usepackage{pgfpages}
\setbeamertemplate{caption}[numbered]
\setbeamertemplate{caption label separator}{: }
\setbeamercolor{caption name}{fg=normal text.fg}
\beamertemplatenavigationsymbolsempty
% remove section numbering
\setbeamertemplate{part page}{
  \centering
  \begin{beamercolorbox}[sep=16pt,center]{part title}
    \usebeamerfont{part title}\insertpart\par
  \end{beamercolorbox}
}
\setbeamertemplate{section page}{
  \centering
  \begin{beamercolorbox}[sep=12pt,center]{section title}
    \usebeamerfont{section title}\insertsection\par
  \end{beamercolorbox}
}
\setbeamertemplate{subsection page}{
  \centering
  \begin{beamercolorbox}[sep=8pt,center]{subsection title}
    \usebeamerfont{subsection title}\insertsubsection\par
  \end{beamercolorbox}
}
% Prevent slide breaks in the middle of a paragraph
\widowpenalties 1 10000
\raggedbottom
\AtBeginPart{
  \frame{\partpage}
}
\AtBeginSection{
  \ifbibliography
  \else
    \frame{\sectionpage}
  \fi
}
\AtBeginSubsection{
  \frame{\subsectionpage}
}
\usepackage{iftex}
\ifPDFTeX
  \usepackage[T1]{fontenc}
  \usepackage[utf8]{inputenc}
  \usepackage{textcomp} % provide euro and other symbols
\else % if luatex or xetex
  \usepackage{unicode-math} % this also loads fontspec
  \defaultfontfeatures{Scale=MatchLowercase}
  \defaultfontfeatures[\rmfamily]{Ligatures=TeX,Scale=1}
\fi
\usepackage{lmodern}
\ifPDFTeX\else
  % xetex/luatex font selection
\fi
% Use upquote if available, for straight quotes in verbatim environments
\IfFileExists{upquote.sty}{\usepackage{upquote}}{}
\IfFileExists{microtype.sty}{% use microtype if available
  \usepackage[]{microtype}
  \UseMicrotypeSet[protrusion]{basicmath} % disable protrusion for tt fonts
}{}
\makeatletter
\@ifundefined{KOMAClassName}{% if non-KOMA class
  \IfFileExists{parskip.sty}{%
    \usepackage{parskip}
  }{% else
    \setlength{\parindent}{0pt}
    \setlength{\parskip}{6pt plus 2pt minus 1pt}}
}{% if KOMA class
  \KOMAoptions{parskip=half}}
\makeatother
\setlength{\emergencystretch}{3em} % prevent overfull lines
\providecommand{\tightlist}{%
  \setlength{\itemsep}{0pt}\setlength{\parskip}{0pt}}
% Slide layout defaults for warning-clean scientific decks.
\usepackage{etoolbox}
\IfFileExists{xurl.sty}{\usepackage{xurl}}{}
\IfFileExists{seqsplit.sty}{\usepackage{seqsplit}}{\newcommand{\seqsplit}[1]{#1}}
\protected\def\breakseq#1{\seqsplit{#1}}
\protected\def\breaktt#1{\begingroup\ttfamily\seqsplit{#1}\endgroup}
\setlength{\emergencystretch}{6em}
\tolerance=5000
\hbadness=10000
\hfuzz=1pt
\setlength{\tabcolsep}{2pt}
\AtBeginEnvironment{longtable}{\tiny\renewcommand{\arraystretch}{0.86}\setlength{\tabcolsep}{1pt}}
\AtBeginEnvironment{tabular}{\tiny\renewcommand{\arraystretch}{0.86}\setlength{\tabcolsep}{1pt}}

% Natbib and cross-reference fallbacks — slides don't load natbib
% or cleveref, but combined-PDF manuscript prose may emit these
% commands. The fallback renders citations as a bracketed key list
% and cross-references as detokenized labels so slides don't fail on
% undefined control sequences or raw underscores. \providecommand is
% a no-op if packages load later.
\providecommand{\citep}[1]{[#1]}
\providecommand{\citet}[1]{#1}
\providecommand{\citealp}[1]{#1}
\providecommand{\citeauthor}[1]{#1}
\providecommand{\citeyear}[1]{#1}
\providecommand{\cref}[1]{\texttt{\detokenize{#1}}}
\providecommand{\Cref}[1]{\texttt{\detokenize{#1}}}
\usepackage{bookmark}
\IfFileExists{xurl.sty}{\usepackage{xurl}}{} % add URL line breaks if available
\urlstyle{same}
% fallback for those not using the hyperref driver hyperxmp:
\makeatletter
\@ifundefined{xmpquote}{\newcommand{\xmpquote}[1]{#1}}{}
\makeatother
\hypersetup{
  hidelinks,
  pdfcreator={LaTeX via pandoc}}

\author{\texorpdfstring{}{}}
\date{}

\begin{document}

\section{Results}\label{sec:results}

\begin{frame}[fragile,allowframebreaks]{Results}
After running \breaktt{scripts/run\_prose\_pipeline.py} with the bundled
\breaktt{manuscript/config.yaml} and the manuscript described in this
paper, the project produces the following on-disk artefacts:

\begin{itemize}
\tightlist
\item
  \breaktt{output/manuscript\_report.json} --- the raw
  \breaktt{ManuscriptReport} JSON.
\item
  \breaktt{output/checks.json} --- one \texttt{CheckResult} per
  configured check.
\item
  \breaktt{output/review\_report.md} --- the assembled markdown review
  report.
\item
  \breaktt{output/figures/section\_word\_counts.png} --- per-file word
  counts.
\item
  \breaktt{output/figures/readability\_metrics.png} --- Flesch /
  Flesch-Kincaid / Gunning Fog per file.
\item
  \breaktt{output/figures/citation\_density.png} --- citations per 1000
  words per file.
\item
  \breaktt{output/data/manuscript\_variables.json} --- substitution
  variables for the abstract.
\end{itemize}

By design the three figures are \textbf{standalone CI/diagnostic
artefacts}, not embedded manuscript images: this is a prose-review
template, so its own rendered PDF deliberately contains no figure floats
and runs cleanly through the very prose gate it documents. A fork that
wants figures inside the manuscript should add real markdown image
references to one of the produced figures (for example
\breaktt{output/figures/section\_word\_counts.png},
\breaktt{output/figures/readability\_metrics.png}, or
\breaktt{output/figures/citation\_density.png}) after those files exist.

Because the pipeline does not consult any external service, every
artefact is reproducible from the manuscript text + \texttt{config.yaml}
alone. A second run on the same inputs produces byte-identical JSON
(modulo timestamp metadata in any caches the project later adds).

The pass/fail status of each configured check is recorded in
\breaktt{output/checks.json}. With the bundled configuration:

\begin{itemize}
\tightlist
\item
  \textbf{\breaktt{grade\_level\_in\_band}} --- the weighted average FKGL
  across this manuscript is reported in the \breaktt{output/checks.json}
  \texttt{details.value} field. The default band \texttt{10.0–18.0} is
  intentionally generous; tighten it for editorial review of
  public-facing material.
\item
  \textbf{\breaktt{citation\_density\_above\_floor}} --- citations per
  1000 words must meet \breaktt{prose.citation\_density\_min\_per\_1000}.
  The bundled config sets the floor to \texttt{0.0} (disabled) so the
  run is green out of the box; researchers should raise it to match the
  publication target.
\item
  \textbf{\breaktt{no\_skipped\_heading\_levels}} --- every file in this
  manuscript uses contiguous heading levels.
\item
  \textbf{\breaktt{every\_file\_has\_h1}} --- every prose file starts
  with an H1.
\item
  \textbf{\breaktt{bibliography\_consistency}} --- every
  \texttt{\breakseq{[@key]}} in the prose resolves against
  \breaktt{manuscript/references.bib}.
\end{itemize}

The figures in \texttt{output/figures/} are colour-blind-safe (Wong 2011
palette) \breakseq{[@wong2011points]}, 300 dpi, and PNG-only for archival
stability. They are referenced in \breakseq{[@sec:pipeline\_internals]} where
we walk through the on-disk artefact set in full.

The full pass/fail summary lands in \breaktt{output/review\_report.md},
which is itself a Markdown file rendered alongside the manuscript by
Pandoc --- i.e.~the project reports on itself in the same compilation
step that produces the manuscript PDF.
\end{frame}

\end{document}
